% ===================================================================
%  LENTELĖ Nr. 10 — Jūra. Uostas.
%  Traduction 2026 d'apres texto/io, controlee sur texto/fr, texto/ru
%  et texto/cs. Consigne : voir texto/lt/10-lentele-01.tex.
%
%  LE TABLEAU 9 A MONTRE UN PAYS SANS MONTAGNE. LE TABLEAU 10 MONTRE
%  UN PAYS QUI A UNE MER, MAIS PAS CELLE-CI. La Lituanie touche la
%  Baltique par quatre-vingt-dix kilometres de cote basse, de dunes
%  et de lagune : ni falaise, ni cap, ni — et c'est le point — DE
%  MAREE. La Baltique n'en a presque pas.
%
%  DE SORTE QUE « potvynis » FAIT DOUBLE EMPLOI. Au tableau 9 le mot
%  dit la crue de la riviere qui a change le chemin en marais ; ici,
%  avec « atoslūgis » son contraire, il doit dire le flux et le
%  reflux de l'ocean. C'est exactement l'operation tcheque du tableau
%  10 de la treizieme colonne — « útes » pour la falaise et pour
%  l'ecueil — mais elle porte ici sur la mer et non sur la montagne :
%  UNE LANGUE PRETE A CE QU'ELLE N'A PAS LES MOTS DE CE QU'ELLE A.
%
%  ET LE MOT D'ECOLE EST UN EMPRUNT LA OU LE MOT DE TOUS LES JOURS
%  N'EN EST PAS UN. « skardis » nomme un talus abrupt, de riviere ou
%  de mer indifferemment ; la falaise de la lecon de geographie se
%  dit « klifas », et l'ecueil « rifas ». Le mot general est du fonds,
%  le terme technique est pris au-dehors. Le tableau 5 avait donne la
%  meme partition sur le chantier, entre le metier qu'on avait et le
%  metier qu'on a recu.
%
%  LES NOMS DE NAVIRES, EUX, SE COMPOSENT, comme se composaient les
%  ustensiles du tableau 6 : « garlaivis » le vapeur, « burlaivis »
%  le voilier, « povandeninis laivas » le sous-marin, « šarvuotis »
%  le cuirasse — de « šarvai » l'armure —, « naikintuvas » le
%  contre-torpilleur — de « naikinti » detruire —, « minininkas » le
%  torpilleur, « vilkikas » le remorqueur — de « vilkti » trainer.
%  Restent etrangers « kreiseris », « fregata », « torpeda »,
%  « dokas », « kranas » : la coque se nomme, la machine s'emprunte.
%
%  ET LE PHARE EST « švyturys », CE QUI BRILLE : le suffixe des
%  ustensiles, une fois encore, et cette fois sur la chose la plus
%  visible du tableau.
%
%  UN GENITIF TENU PAR UNE APPOSITION : le pont (48) d'un voilier
%  (49). Range a la lituanienne le voilier passerait devant son pont ;
%  deux datifs juxtaposes gardent l'ordre de la source.
% ===================================================================

%%K t10-apar-1 apar
\VUsaut{4.0mm}
\VUtitre{15.0pt}{60}{LENTELĖ Nr. 10}
\VUsaut{3.0mm}
\VUpk{11.4pt}{40}{Jūra. --- Uostas.}
\VUsaut{3.0mm}

%%K t10-01-1 p
1. --- Vasarą, kol vieni ieško ramybės ir vėsos kalnuose, kiti mieliau vyksta į
pajūrį. Taip padarėme pernai, nes labai mėgstame \VUgras{Vandenyną}
\textsuperscript{(1)}.

%%K t10-02-1 p
2. --- Man didelis malonumas žvelgti į tą milžinišką vandens platybę, kuri driekiasi
iki \VUgras{horizonto} \textsuperscript{(2)}, ir matyti išplaukiančius į ilgą
\VUgras{kelionę} milžiniškus \VUgras{garlaivius} \textsuperscript{(3)}, į kuriuos
sėda keleiviai, vykstantys į Ameriką.

%%K t10-03-1 p
3. --- Todėl mano tėvas, kuris žino mano polinkį, atostogoms visada pasirenka kokį
labai ramų paplūdimį netoli vieno iš mūsų didžiųjų \VUgras{karinių uostų.}

%%K t10-03-2 p
Per paskutinę mūsų viešnagę \VUgras{eskadra} atplaukė stoti už milžiniško
\VUgras{molo} \textsuperscript{(4)}, kuris uždaro ir saugo \VUgras{karinį uostą}
\textsuperscript{(5)}, ir man pavyko, kaip ir norėjau, pasigėrėti įvairiais
\VUgras{kariniais laivais,} mažu \VUgras{povandeniniu laivu} \textsuperscript{(10)},
kuris neria ir dingsta po vandeniu, ir taip pat milžinišku \VUgras{šarvuočiu}
\textsuperscript{(9)}, kuris iki šiol bijojo bemaž vien \VUgras{minininko}
\textsuperscript{(8)} atakų.

%%K t10-tit-2 sub
\VUsaut{3.0mm}
\VUpk{10.2pt}{40}{Mažieji laivai.}
\VUsaut{3.0mm}

%%K t10-04-1 p
4. --- Šis jau mokėjo labai \VUgras{vikriai} rasti storo metalinio šarvo silpnąją
vietą, kad įsmeigtų ten pavojingą \VUgras{torpedą,} kurios sprogimas sukelia laivo
žūtį; bet vis dėlto jo reikėjo bijoti mažiau negu povandeninio laivo, nes naikintuvai
lengvai jį persekioja.

%%K t10-05-1 p
5. --- Galėjau pamatyti ir greitus šarvuotus \VUgras{kreiserius}
\textsuperscript{(6)}, beveik tokius pat nepažeidžiamus kaip tikras šarvuotis, kelis
\VUgras{transporto laivus} \textsuperscript{(12)}, tris ar keturis \VUgras{patrankinius
laivus} \textsuperscript{(7)} ir galiausiai vieną \VUgras{fregatą}
\textsuperscript{(11)}. \VUgras{To laivyno reginys, kai jis manevruoja keldamas
inkarą, yra įdomiausias.}

%%K t10-06-1 p
6. --- Koks statybos skirtumas tarp tų didelių plieno krūvų ir elegantiškų
\VUgras{trimasčių} ar grakščių \VUgras{burlaivių} \textsuperscript{(13)}, dar ir
dabar vartojamų prekybos laivyne, kuriuos mačiau įplaukiančius į uostą visomis burėmis,
kai vaikščiojau \VUgras{prieplauka} \textsuperscript{(14)}. Tiesa, šie prarado daug
savo pranašumų, \VUgras{kai} \VUgras{vėjas} buvo priešingas. Jiems teko nuleisti
visas bures, kad grįžtų į baseiną, ir \VUgras{vilkiko} \textsuperscript{(16)} reikėjo,
kad juos pasiimtų ir atvestų, kaip paprastas \VUgras{baržas} \textsuperscript{(17)},
prie krantinės.

%%K t10-tit-3 sub
\VUsaut{3.0mm}
\VUpk{10.2pt}{40}{Prekybos uostas.}
\VUsaut{3.0mm}

%%K t10-07-1 p
7. --- Kas įdomu uoste, apie kurį kalbu, yra tai, kad jame yra ne tik karinis uostas,
dirbtinai įrengtas milžiniškais darbais, bet ir nuostabus \VUgras{prekybos uostas,}
esantis didelės upės \VUgras{žiotyse.}

%%K t10-08-1 p
8. --- Žemės liežuvis, skiriantis abu baseinus, sutvirtintas ir ginamas
\VUgras{baterijų} ir \VUgras{bastionų} eile, kurie išsikiša į jūrą. Reikia
specialaus leidimo pasivaikščioti \VUgras{pylimais} \textsuperscript{(15)}. Lengvai jį
gavau per dėdę, kuris yra \VUgras{laivų statyklų} \textsuperscript{(18)} inžinierius,
ir nuėjau apžiūrėti laivų, kurie statomi po plačiais angarais.

%%K t10-09-1 p
9. --- Net dalyvavau šarvuočio nuleidime ir prisimenu jaudinančią akimirką, kurią
pajuto visa minia, kai milžiniška masė, palikta pati sau, ėmė slysti savo lopšiniu
rėmu ir atplaukė plūduriuoti upės vandeniu.

%%K t10-10-1 p
10. --- Norint pasiekti laivų statyklas, einama per \VUgras{pratybų lauką}
\textsuperscript{(19)}, kur garnizono kareiviai kasdien ateina mankštintis, ir per
geležinį \VUgras{tiltelį} \textsuperscript{(20)}, kuris jungia paradų aikštę su
\VUgras{arsenalu} \textsuperscript{(21)}.

%%K t10-tit-4 sub
\VUsaut{3.0mm}
\VUpk{10.2pt}{40}{Arsenalas ir dokai.}
\VUsaut{3.0mm}

%%K t10-11-1 p
11. --- Su dėde aplankiau tą didelį statinį, pilną \VUgras{patrankų}
\textsuperscript{(22)}, \VUgras{sviedinių} \textsuperscript{(23)} ir \VUgras{haubicų
sviedinių.} Taip pat mačiau \VUgras{metalo liejyklą} \textsuperscript{(24)}, kurioje
gaminami garlaivių \VUgras{katilai} \textsuperscript{(25)}.

%%K t10-12-1 p
12. --- Visai netoli yra \VUgras{dokai} \textsuperscript{(26)} --- remonto dokai --- ir
labai žymūs \VUgras{plaukiojantys dokai} \textsuperscript{(27)}, kuriuose galėjau
matyti visai iš arti, ant taisytino laivo, laivo \VUgras{sraigtą}
\textsuperscript{(28)}, \VUgras{kilį} \textsuperscript{(29)} ir \VUgras{vairą}
\textsuperscript{(30)}.

%%K t10-13-1 p
13. --- Prekybos uostas lygiai taip pat vertas studijų kaip karinis uostas, ir
praleidau daug valandų žioplinėdamas ant krantinių, kur pririšti \VUgras{irklaračiai
laivai} \textsuperscript{(31)}, kurie plaukia upe atgal, ir žiūrėdamas, kaip veikia
\VUgras{stiebiniai kranai} \textsuperscript{(32)}, kurie kelia kaip plunksneles
milžiniškus krovinius.

%%K t10-tit-5 sub
\VUsaut{4.0mm}
\VUpk{10.2pt}{40}{Locmanas.}
\VUsaut{3.0mm}

%%K t10-14-1 p
14. --- Gėrėjausi \VUgras{transatlantiniais laivais} \textsuperscript{(34)},
išplaukiančiais iš reido, su \VUgras{vėliavomis} \textsuperscript{(35)} plazdančiomis
ore, vedamais sumanaus \VUgras{locmano} \textsuperscript{(36)}, kuris nuostabiai moka
laikytis \VUgras{farvaterio,} ženklinamo \VUgras{plūdurais} \textsuperscript{(40)} ir
\VUgras{ženklais,} vengdamas \VUgras{baidokų} \textsuperscript{(41)}, kurie sunkiai
valdomi savo vienintele keturkampe bure. Įžiūrėjau \VUgras{priešakyje}
\textsuperscript{(37)} --- laivapriekyje --- antrosios klasės \VUgras{kajutes}
\textsuperscript{(39)}, o \VUgras{laivagalyje} \textsuperscript{(38)} --- ties
laivagaliu --- pirmosios klasės salonus.

%%K t10-15-1 p
15. --- Kartais taip pat mačiau, kaip kraunamas \VUgras{krovininis laivas}
\textsuperscript{(42)}, į kurį ką tik buvo sukrautos prekės iš \VUgras{sandėlio}
\textsuperscript{(43)} ir iš \VUgras{jūrų stoties} \textsuperscript{(44)}, atgabentos
daugybe \VUgras{krantinės geležinkelio} \textsuperscript{(45)} traukinių.

%%K t10-tit-6 sub
\VUsaut{3.0mm}
\VUpk{10.2pt}{40}{Irklavimo pramoga.}
\VUsaut{3.0mm}

%%K t10-16-1 p
16. --- Dažnai irstydavausi \VUgras{uosto baseine} \textsuperscript{(53)}, į kurį
ateina prisiglausti \VUgras{valtys} \textsuperscript{(46)}, ir man buvo džiaugsmas
valdyti \VUgras{irklą} \textsuperscript{(47)}. Užlipau ant \VUgras{denio}
\textsuperscript{(48)}, ant vienos \VUgras{burvaltės} \textsuperscript{(49)}, lipau
\VUgras{lynais} \textsuperscript{(52)} ant \VUgras{stiebo} \textsuperscript{(51)} į
\VUgras{rėjas} \textsuperscript{(50)}, paskui pakartojau savo pasiplaukiojimą
\VUgras{jole} \textsuperscript{(54)} tarp sunkių \VUgras{žvejų valčių}
\textsuperscript{(55)}, kur džiūsta \VUgras{tinklai} \textsuperscript{(56)}.

%%K t10-17-1 p
17. --- Ant \VUgras{krantinės} \textsuperscript{(58)} pro mane praėjo pačios
įvairiausios \VUgras{prekės} \textsuperscript{(59)}, sudėtos \VUgras{dėžėse}
\textsuperscript{(60)} arba \VUgras{ryšuliuose} \textsuperscript{(61)}. Jos sudarė
baržų \VUgras{krovinį} \textsuperscript{(63)}, ir galingi \VUgras{kranai}
\textsuperscript{(62)} juos nukėlė ir padėjo ant \VUgras{sunkvežimių}
\textsuperscript{(64)}, o \VUgras{vežėjai} juos deklaravo ir sumokėjo mokesčius
\VUgras{muitinėje} \textsuperscript{(65)}.

%%K t10-18-1 p
18. --- Galiausiai, kai priėjau \VUgras{pasukamąjį tiltą} \textsuperscript{(66)} ar
\VUgras{šliuzą} \textsuperscript{(69)}, eidamas pro \VUgras{žemsiurbę}
\textsuperscript{(68)}, kuri valo \VUgras{kanalą} \textsuperscript{(67)}, išlipau į
prieplauką prie \VUgras{semaforo} \textsuperscript{(70)}.

%%K t10-19-1 p
19. --- Kitoje tos prieplaukos pusėje kaip tik priplaukia \VUgras{laivavaltės}
\textsuperscript{(71)}, kurios veža į krantą eskadros karininkus. Nuostabus reginys ---
matyti, kaip jūreiviai ritmingai kelia savo irklus, kol valtis, nešama \VUgras{bangos}
\textsuperscript{(72)}, lengvai priplaukia prie išlaipinimo laiptų. Toje vietoje
geriau galima stebėti \VUgras{jūros kaitą} ir \VUgras{potvynio} bei
\VUgras{atoslūgio} judėjimą \VUgras{paplūdimyje} \textsuperscript{(73)}.

%%K t10-tit-7 sub
\VUsaut{3.0mm}
\VUpk{10.2pt}{40}{Klubas.}
\VUsaut{3.0mm}

%%K t10-20-1 p
20. --- Grįžti namo pasinaudojau \VUgras{elektriniu tramvajumi}
\textsuperscript{(74)}, kurio \VUgras{stotelė} \textsuperscript{(75)} yra prie
\VUgras{stulpo,} pastatyto prieš karininkų klubą.

%%K t10-20-2 p
Nuo klubo \VUgras{balkono} \textsuperscript{(76)}, papuošto \VUgras{inkarais}
\textsuperscript{(77)}, patrankomis ir sviediniais, dažnai mačiau puikų jūrų
karininkų būrį: \VUgras{leitenantus} \textsuperscript{(78)} su savo špagomis,
\VUgras{artilerijos} \textsuperscript{(79)} ir \VUgras{pėstininkų}
\textsuperscript{(80)} \VUgras{kapitonus} su savo \VUgras{kardais.} Už jų buvo
\VUgras{jūreivis} \textsuperscript{(81)}, kuris paduodavo jiems \VUgras{žiūroną,} kai
jie norėdavo įžiūrėti, kas vyksta toli.

%%K t10-21-1 p
21. --- Mačiau ir jaunus \VUgras{jaunesniuosius leitenantus} \textsuperscript{(82)},
gaunančius nurodymus iš savo \VUgras{vado} \textsuperscript{(83)} ar iš
\VUgras{admirolo} \textsuperscript{(84)}, kol \VUgras{bocmanas}
\textsuperscript{(85)} laukė, kad nuneštų juos į laivą.

%%K t10-22-1 p
22. --- Nuo to balkono vaizdas puikus: matyti visas paplūdimys su savo
\VUgras{palapinėmis} \textsuperscript{(86)} ir \VUgras{kabinomis}
\textsuperscript{(87)}, ir matyti, maudymosi valandą, \VUgras{besimaudantieji}
\textsuperscript{(88)}, kurie džiaugsmingai išdykauja prieš \VUgras{kazino}
\textsuperscript{(89)}.

%%K t10-23-1 p
23. --- Paplūdimio gale krantas sudaro mažą \VUgras{uolėtą} \VUgras{pusiasalį}
\textsuperscript{(91)}, į kurį atsimuša \VUgras{didžioji banga} \textsuperscript{(92)}
ir ant kurio pastatytas \VUgras{švyturys} \textsuperscript{(93)}, naktį rodantis kelią
plaukiantiesiems. Tas pusiasalis su žeme susisiekia tik labai siaura \VUgras{sąsmauka}
\textsuperscript{(90)}.

%%K t10-tit-8 sub
\VUsaut{3.0mm}
\VUpk{10.2pt}{40}{Bendras pakrančių vaizdas.}
\VUsaut{3.0mm}

%%K t10-24-1 p
24. --- \VUgras{Klifas} \textsuperscript{(94)} tęsiasi, perskeltas jūros, kuri
įnoringai išraižo jame \VUgras{įlankų} \textsuperscript{(95)} ir
\VUgras{iškyšulių} (ragų) eilę, padarydama jį vaizdingiausią.

%%K t10-24-2 p
Ant \VUgras{plokščiakalnio} \textsuperscript{(96)}, kuris dominuoja virš
\VUgras{sąsiaurio} \textsuperscript{(98)}, yra labai svarbus \VUgras{fortas}
\textsuperscript{(97)} ir kelios „vilos“.

%%K t10-25-1 p
25. --- Tas sąsiauris atskiria nuo žemyno labai pavojingą \VUgras{salą}
\textsuperscript{(99)}, apsuptą \VUgras{rifų} \textsuperscript{(100)} ir
\VUgras{seklumų,} kur valtys ir net dideli laivai kartais užplaukia ant seklumos.
Nepaisant pagalbos skubumo ir greito \VUgras{gelbėjimo valčių} atplaukimo, tos
\VUgras{laivų žūtys} baisios, ir per lygiadienio \VUgras{audras} keli laivai žuvo su
žmonėmis ir prekėmis.

%%K t10-25-2 p
Nepaisant tos kaimynystės, \VUgras{uostas} labai lankomas jūreivių.

\VUsaut{6.0mm}
\VUfilet{22mm}
